\section{Why Predicting Band Power Is Efficient}\label{sec:theory}
We motivate the target under an idealized model; the statements and proofs are in
Appendix~\ref{app:theory}. Before the logarithm and numerical floor, the main target,
Eq.~\eqref{eq:target}, is a per-band constant times the band-averaged periodogram. The exact
log-noise identities below apply to the zero-floor idealization with independent Fourier bins
and exactly flat in-band spectra. The implemented $\log(P+\varepsilon)$ and real EEG need not
satisfy them. The analysis motivates the target and suggests empirical hypotheses; it is not a
proof of optimality or an optimization-speed guarantee.

\noindent\emph{Model.} Given a brain state $\theta$, a patch is a segment of a zero-mean
stationary Gaussian process with spectral density $S_\theta$; the state acts on the data only
through $S_\theta$. Asymptotically, the discrete Fourier coefficients at distinct frequencies
are independent circular complex Gaussians, so the periodogram ordinates are $I_k\approx
S_\theta(f_k)E_k$ with $E_k\sim\mathrm{Exp}(1)$, and the phases are uniform and independent of
$\theta$ \cite{brillinger1981}. The Whittle likelihood \cite{whittle1953} depends on the data
only through $\{I_k\}$.

\noindent\emph{What band power keeps.} (i) In the Whittle model, the periodogram is sufficient
for its spectral parameters and the marginal phases are ancillary (Prop.~\ref{prop:suff});
this does not establish independence from visible context. (ii) With exactly constant in-band
spectra and observation restricted to those bins, the band-averaged periodogram is sufficient
for the band levels and is their maximum-likelihood estimate (Prop.~\ref{prop:mle}). (iii)
For its zero-floor logarithm, $\Var[\log P_b\mid\theta]=\psi'(m_b)\approx 1/m_b$ is independent
of the state (Prop.~\ref{prop:var}). This motivates variance stabilization, but neither
sufficiency for arbitrary EEG labels nor an exact noise model for Eq.~\eqref{eq:target} follows.

\noindent\emph{A possible learning advantage.} Under squared error, prediction risk splits
into error relative to $g(\bm{c})=\E[Y\mid\bm{c}]$ and the Bayes risk given the fixed context
(Prop.~\ref{prop:decomp}). Only the former can be reduced by improving the predictor.
Under an additional independent-patch idealization, the zero-floor band-power predictor is
$\E[\log s_b(\theta)\mid\bm{c}]+\text{const}$, a posterior average of spectral levels,
rather than the state itself. Its conditional variance contains both $\psi'(m_b)$ and posterior
uncertainty about $\log s_b(\theta)$. In this same independence model the zero-mean waveform
predictor would be zero. A separate correlated Gaussian model instead yields Wiener
interpolation conditional on the state, with additional uncertainty when the state is unknown.
For either target, the full conditional covariance contributes Jacobian-weighted gradient
noise (Prop.~\ref{prop:grad}). Standard SGD upper bounds can worsen with gradient variance
\cite{ghadimi2013}, but the two targets also differ in conditional mean, dimension, scale and
Jacobian. These identities suggest a possible advantage of a compact spectral target; they do
not establish an ordering of gradient noise or steps to a given downstream score.

\noindent\emph{Predictions.} (P1) At matched backbone, data and steps, band power should give
representations that are easier to read out linearly, especially with few labels. (P2) It should
reach a given quality with fewer pretraining steps and smaller models. (P3) Reconstructing
band-filtered \emph{waveforms} need not yield the same gains, since they still require sample-level prediction.
(P4) The advantage should shrink for tasks that depend on transient waveform shape or
cross-channel phase relations. Sec.~\ref{sec:results} tests these.
